\documentclass[11pt,a4paper]{article}
\usepackage[margin=2.5cm]{geometry}
\usepackage{amsmath,amssymb,bm}
\usepackage{booktabs}
\usepackage[colorlinks=true,linkcolor=blue,citecolor=blue,urlcolor=blue]{hyperref}

\newcommand{\vx}{\bm{x}}
\newcommand{\E}{\mathbb{E}}
\newcommand{\Ev}{E^{\mathrm{v}}}
\newcommand{\grad}{\nabla}

\title{Physical diffusion models with fast latent units:\\ model, local training and measured statistics}
\author{Working note, branch \texttt{latent\_physical\_SGM}}
\date{9 November 2026}

\begin{document}
\maketitle

\begin{abstract}
We extend the physical diffusion model --- a lattice of nonlinear (Duffing) oscillators, one per pixel, whose
energy is trained per time slice by score matching --- with latent degrees of freedom: hidden units that couple to
the pixels and equilibrate much faster than they do. Integrated out, the hidden units give the pixels a free
energy with a global, many-body part, in particular a ``which prototype'' variable that a two-body pixel energy
lacks. We derive the free energy for Ising, Potts and per-pixel hidden units, the score-matching objective and its
gradient, and show that the gradient is a local three-factor rule. On a device the hidden statistics must be
measured from a few samples; we derive the bias this causes, remove it with a first- and second-order jackknife,
and describe an optimiser (the visit-count rule) that keeps rarely visited hidden states alive. We close with the
reverse (generative) process, the physical conditions for the fast-hidden idealisation, the evaluation, and the
main numerical results so far.
\end{abstract}

\tableofcontents

%======================================================================================================================
\section{Setting: the physical diffusion model}

\paragraph{Forward process.} The data $\vx\in\mathbb R^N$ (pixels, $N=784$ at $28\times28$, $N=196$ at $14\times14$,
normalised) are diffused by an Ornstein--Uhlenbeck process at temperature $T$,
\begin{equation}
  d\vx = -\vx\,dt + \sqrt{2T}\,d\bm W_t ,\qquad T=0.02,\qquad t\in[0,5],
\end{equation}
whose stationary density is $\propto e^{-\|\vx\|^2/2T}$. The marginal at time $t$ is $p_t$.

\paragraph{Time slices.} The interval $t\in[0,5]$ is cut into 100 slices. Every parameter of the model has its own
value in each slice; slice $k$ is trained on fresh samples of $p_{t_k}$ and warm-started from slice $k-1$. During
generation the parameters are interpolated between the slice knots.

\paragraph{Energy-based score.} At each slice the model density is $p_\theta(\vx)\propto e^{-F_\theta(\vx)/T}$, so
its score is $\grad\log p_\theta=-\grad F_\theta/T$. Without hidden units $F_\theta=\Ev_\theta$, the visible energy.

\paragraph{Visible energy.} One oscillator per pixel, with self terms and pair couplings,
\begin{equation}
  \Ev(\vx)=\sum_i\Big[\tfrac{k^{\rm lin}_i}{2}x_i^2+\tfrac{k^{\rm duff}_i}{4}x_i^4+\tfrac{k^{(6)}_i}{6}x_i^6+\beta_ix_i\Big]
  +\sum_{(i,j)\in\mathcal P}\Big[c^{\rm lin}_{ij}(x_i-x_j)^2+c^{\rm om}_{ij}\,x_i^2x_j+\tfrac{c^{\rm duff}_{ij}}{4}(x_i-x_j)^4\Big].
  \label{eq:Ev}
\end{equation}
It is linear in its parameters $\theta_v$. The paper's model (``power 6'') keeps $k^{(6)}_i\ge0.01$ for
confinement; ``power 4'' sets $k^{(6)}=0$ and confines with $k^{\rm duff}_i\ge0.01$, $c^{\rm duff}\ge0$. A pixel's
self energy is a double well wherever $k^{\rm lin}_i<0$. The pair set $\mathcal P$ contains all pairs within a
Chebyshev distance (14 at $28\times28$: 180{,}810 pairs, 59\% of all pairs; 7 at $14\times14$). The three pair terms
are the linear spring ($c^{\rm lin}$), an asymmetric ``optomechanical'' term ($c^{\rm om}$) and a quartic pair term
($c^{\rm duff}$). Options \texttt{--radius}, \texttt{--neighbourhood von\_neumann} and \texttt{--pair-terms} restrict
$\mathcal P$ and the pair terms (Section~\ref{sec:results}).

%======================================================================================================================
\section{Latent degrees of freedom}

\subsection{Joint energy and marginal}
Add hidden units $\bm y\in\mathbb R^K$ that couple bilinearly to the pixels and not to each other (a bipartite,
RBM-like structure):
\begin{equation}
  E_\theta(\vx,\bm y)=\Ev(\vx)+\sum_{j=1}^K\big[U_j(y_j)-y_j\,a_j(\vx)\big],\qquad a_j(\vx)=b_j+\bm W_j\cdot\vx ,
\end{equation}
with $U_j$ the unit's own potential, e.g.\ a double well. The joint energy is linear in all parameters. Integrating
out the hidden units gives the marginal $p_\theta(\vx)\propto\int d\bm y\,e^{-E_\theta/T}=e^{-F_\theta(\vx)/T}$ with
the free energy
\begin{equation}
  F(\vx)=\Ev(\vx)-T\sum_j\log Z_j\big(a_j(\vx)\big),\qquad Z_j(a)=\int dy\,e^{-(U_j(y)-ya)/T}.
\end{equation}
$T\log Z_j$ is the cumulant generating function of $y_j$; writing $m_j$, $v_j$, $k_{3,j}$ for the conditional mean,
variance and third cumulant of $y_j$ given $\vx$,
\begin{equation}
  T\partial_a\log Z_j=m_j,\qquad\partial_a m_j=v_j/T,\qquad\partial_a v_j=k_{3,j}/T .
\end{equation}

\paragraph{Idealisation.} We always integrate the hidden units out exactly. Physically this is the limit of
hidden units that equilibrate infinitely fast compared with the pixels (Section~\ref{sec:physical}).

\subsection{Ising units}
For a deep double well with minima at $\pm y_0$, $y_j\to y_0h_j$ with $h_j=\pm1$ ($y_0$ absorbed into $\bm W,b$):
\begin{equation}
  Z_j=2\cosh(a_j/T),\qquad m_j=\tanh(a_j/T),\qquad v_j=1-m_j^2,\qquad k_{3,j}=-2m_jv_j .
\end{equation}
With $K$ Ising units the model is a mixture of $2^K$ exponentially tilted copies of the visible model,
\begin{equation}
  p_\theta(\vx)\propto\sum_{\bm h\in\{\pm1\}^K}\exp\Big(-\big[\Ev(\vx)-\textstyle\sum_jh_ja_j(\vx)\big]/T\Big).
\end{equation}
One unit with $\bm W_1$ along the Fisher (LDA) direction separates two classes; $P(h=+1|\vx)=\sigma(2a_1/T)$ is then the
LDA posterior. So far only $K=1$ has been trained; several units need a symmetry-breaking initialisation, because
$\bm W=0$ is a degenerate stationary point (Section~\ref{sec:training}).

\subsection{Potts units}
A Potts unit is \emph{one} hidden unit with $M$ mutually exclusive states $m=1\dots M$ (exactly one active). State
$m$ has its own couplings to all pixels and its own field:
\begin{equation}
  F(\vx)=\Ev(\vx)-T\log\sum_{m=1}^Me^{a_m(\vx)/T},\qquad p_m(\vx)=\frac{e^{a_m/T}}{\sum_ke^{a_k/T}} .
\end{equation}
The model is a mixture of $M$ tilted copies of the visible model; each state acts as a prototype. ``Potts 128'' is
one unit with 128 states, not 128 units. In hardware a Potts unit is an $M$-well oscillator, or $M$ bistable units
with strong lateral inhibition (winner-take-all); the code idealises it as exactly one-hot.

\subsection{Spring coupling and per-pixel units}
\paragraph{Spring coupling (\texttt{--quad}).} The input field also contains the squared pixels,
\begin{equation}
  a_m(\vx)=b_m+\bm W_m\cdot\vx+\bm Q_m\cdot\vx^{2},\qquad \bm G_m\equiv\grad a_m=\bm W_m+2\bm Q_m\circ\vx .
\end{equation}
State $m$ then adds $-Q_{mi}x_i^2$ to pixel $i$'s energy: it changes the pixels' stiffness, so the mixture
components can differ in width and shape, not only in position. A bilinear coupling alone only shifts the
components (for a quadratic $\Ev$ exactly), so class-specific shape would otherwise have to come from the
anharmonic part of $\Ev$.

\paragraph{Per-pixel units (\texttt{--local}).} One Ising unit per pixel, coupled only to its own pixel,
$a_i=c_i+w_ix_i$, adds $-T\log2\cosh\big((c_i+w_ix_i)/T\big)$: an asymmetric double well per pixel (ink or
background). ``Potts $M$ + both'' denotes one Potts unit with spring coupling plus the per-pixel units:
\begin{equation}
  F=\Ev-T\log\sum_me^{a_m/T}-T\sum_i\log2\cosh\frac{c_i+w_ix_i}{T}.
  \label{eq:Fboth}
\end{equation}

\subsection{What the latent units add}
\paragraph{Mean force and softening.} With $\bm\psi\equiv\grad F$ (the force a clamp must supply to hold the
pixels), for Ising units
\begin{equation}
  \bm\psi=\grad\Ev-\sum_jm_j\bm W_j,\qquad
  \grad\grad F=\grad\grad\Ev-\frac1T\sum_jv_j\bm W_j\bm W_j^{\!\top},\qquad
  \Delta F=\Delta\Ev-\frac1T\sum_jv_j\|\bm W_j\|^2 .
\end{equation}
The hidden units exert the mean force $\sum_jm_j\bm W_j$ and soften the landscape along $\bm W_j$; near $a_j=0$ a unit
turns a single well into a double well along $\bm W_j$ where $\|\bm W_j\|^2/T$ exceeds the visible curvature.

\paragraph{A many-body, global variable.} The mode-selection force $W_{1i}\tanh(a_1(\vx)/T)$ on pixel $i$ is a
nonlinear function of a global projection of all pixels. A two-body energy cannot represent it exactly: its
forces satisfy $\partial_k\partial_l f_i=0$ for distinct $k,l\neq i$, while the tanh force does not. Expanding,
$-T\log\cosh(a/T)=-a^2/2T+a^4/12T^3-\dots$ contains products of four distinct pixels. The $a^2$ part is itself a pair
term; the hidden units add the non-Gaussian many-body remainder. Empirically, adding more pair couplings to the
paper's model (all pixel pairs) does not reduce 0/1 blends, while one hidden unit does.

\paragraph{A linear hidden mode is not enough.} A harmonic hidden mode $U(y)=\tfrac{k}{2}y^2$ integrates out to
$-a(\vx)^2/2k$: a rank-one softening, i.e.\ another pair term, never the $\pm\bm W$ mixture. The hidden unit must be
nonlinear (bistable or multi-well).

%======================================================================================================================
\section{Training: score matching per slice}\label{sec:training}

\subsection{Objective}
For $p\propto e^{-F/T}$ the implicit score-matching (Hyv\"arinen) objective of slice $t$ is
\begin{equation}
  J_T(\theta)=\E_{\vx\sim p_t}\Big[-\frac{\Delta F}{T}+\frac{\|\grad F\|^2}{2T^2}\Big],\qquad
  \partial_\theta J_T=\E\Big[-\frac{\partial_\theta\Delta F}{T}+\frac{(\partial_\theta\grad F)\cdot\bm\psi}{T^2}\Big].
\end{equation}
It is minimised per slice with Adam on minibatches of 512 fresh samples of $p_t$; the slice result is an
exponential moving average of the iterates. Because $F$ is nonlinear in the hidden parameters, $J_T$ is no longer
a convex quadratic; it has gauge symmetries (sign flips, permutations of identical units), and $\bm W=0$ is a
degenerate stationary point (to second order in $\bm W$ the hidden free energy only adds terms that the visible
bias and pair couplings absorb). Potts states are therefore initialised from $k$-means prototypes (or found from a
random start, which also works), and the single Ising unit along the Fisher direction.

\subsection{Visible parameters}
For $\Ev=\sum_k\theta_k\phi_k$ the form is unchanged, with $\bm\psi$ now including the hidden mean force:
\begin{equation}
  \partial_{\theta_k}J_T=\E\Big[-\frac{\Delta\phi_k}{T}+\frac{\grad\phi_k\cdot\bm\psi}{T^2}\Big].
\end{equation}

\subsection{Ising couplings: a local three-factor rule}
\begin{align}
  \partial_{W_{ji}}J_T&=\frac1{T^2}\E\Big[-m_j\psi_i+2v_jW_{ji}+\frac{x_i}{T}\big(k_{3,j}\|\bm W_j\|^2-v_j\,\bm W_j\cdot\bm\psi\big)\Big],\\
  \partial_{b_j}J_T&=\frac1{T^3}\E\Big[k_{3,j}\|\bm W_j\|^2-v_j\,\bm W_j\cdot\bm\psi\Big].
\end{align}
Each update needs, at pixel $i$, its position $x_i$ and clamp force $\psi_i$; at hidden unit $j$, the statistics of
its own fluctuations ($m_j$; $v_j,k_{3,j}$ follow for Ising units), $\|\bm W_j\|^2$ and the projection
$\bm W_j\cdot\bm\psi$; and the coupling itself. The projection is the change of $a_j$ when the clamped pixels are
nudged along $\bm\psi$; because $a_j$ is linear (or quadratic, with $\bm Q$) in $\vx$, a symmetric pair of nudges
$\pm\varepsilon\bm\psi$ gives it exactly. No free-running negative phase is needed.

\subsection{Potts couplings with spring coupling}
For the model \eqref{eq:Fboth}, per clamped sample define for state $m$: $p_m$, $\bm G_m$, $s_m=\sum_iQ_{mi}$; for
the unit: $\bar{\bm G}=\sum_kp_k\bm G_k$, $V=\sum_kp_k\|\bm G_k-\bar{\bm G}\|^2$, $\bar s=\sum_kp_ks_k$; and the scalar
\begin{equation}
  c_m=\|\bm G_m-\bar{\bm G}\|^2-V+2T(s_m-\bar s)-(\bm G_m-\bar{\bm G})\cdot\bm\psi .
\end{equation}
Then
\begin{align}
  \partial_{b_m}J_T&=\frac{\langle p_mc_m\rangle}{T^3},\qquad
  \partial_{W_{mi}}J_T=\Big\langle p_m\Big[\frac{x_ic_m}{T^3}+\frac{2(G_{mi}-\bar G_i)-\psi_i}{T^2}\Big]\Big\rangle,\\
  \partial_{Q_{mi}}J_T&=\Big\langle p_m\Big[\frac{x_i^2c_m}{T^3}+\frac{2x_i\big(2(G_{mi}-\bar G_i)-\psi_i\big)}{T^2}+\frac2T\Big]\Big\rangle .
\end{align}
Coupler $(m,i)$ needs a state-level scalar ($p_m$, $c_m$) times a pixel signal ($x_i$, $\psi_i$, $\bar G_i$): a
three-factor rule. The per-pixel units follow the Ising rule with a single coupling. This local rule equals the
automatic gradient of the training loss to relative differences of $10^{-16}$--$4\cdot10^{-11}$ on trained models:
every exactly trained model is trained by it, given exact hidden statistics.

\subsection{Contrastive-divergence form}
Hyv\"arinen's relation between one-step contrastive divergence (CD-1) with a short Langevin release and score
matching gives $\langle\partial_\theta F(\vx)\rangle-\langle\partial_\theta F(\vx')\rangle=\tfrac{T^2dt}{2}\partial_\theta J_T$,
i.e.\ a Hebbian term minus the same term after a short release. It trains as well as the exact gradient with exact
hidden averages, but with sampled hidden states the difference of two independent readouts drowns the $O(dt)$
signal; force matching is therefore the rule for a device.

%======================================================================================================================
\section{Training on a device: measured hidden statistics}

\subsection{Plug-in bias}
A device cannot compute $p_m$ or $m_j$; it reads the hidden units $S$ times while the pixels are clamped and
averages. Replacing the exact statistics by these sample means $\hat p_m$ (``measured statistics'', $S$ samples or
\emph{reads} per configuration) gives an unbiased estimate of everything linear in them, but the loss is quadratic
and its parameter gradient at most cubic in the sample means of one configuration. Hence, over $n$ draws,
\begin{equation}
  \E\big[\hat L_n\big]=a+\frac bn+\frac c{n^2},\qquad\text{in particular}\quad \E\big[\hat J_S\big]=J_T-\frac{V}{2ST^2}
\end{equation}
for the loss, with $V$ the variance of the hidden force: the plug-in estimate is biased at order $1/S$.

\subsection{Jackknife corrections}
Split the $S$ draws into halves and quarters and combine the estimates from the same draws:
\begin{align}
  \text{first order (jk):}\quad&2\hat L_S-\tfrac12\big(\hat L_{S/2}^{(1)}+\hat L_{S/2}^{(2)}\big)
     \ \Rightarrow\ \E=a-\frac{2c}{S^2},\\
  \text{second order (jk2):}\quad&\tfrac83\hat L_S-2\,\overline{\hat L_{S/2}}+\tfrac13\,\overline{\hat L_{S/4}}
     \ \Rightarrow\ \E=a .
\end{align}
(The coefficients of $b/n$ cancel as $\tfrac83-2\cdot2+\tfrac13\cdot4=0$, those of $c/n^2$ as
$\tfrac83-2\cdot4+\tfrac13\cdot16=0$.) To first order in the fluctuations neither adds noise. On a device this means
averaging the hidden readouts over one full, two half and four quarter windows. The gradient of the visible
parameters is linear in the sample means and the same for every order.

\subsection{Optimisers and the loss of hidden states}
With few samples, a rarely occupied Potts state is rarely visited: its gradient is zero in most steps and a spike
in the others (unbiased with the jackknife). The optimiser decides whether such a state survives:
\begin{itemize}
\item \textbf{Per-element Adam} divides each coupling's gradient by its own running root-mean-square, so a rarely
  visited state moves at most about $\sqrt\lambda$ as fast as with exact statistics ($\lambda$ = visits per step).
  States below about one visit per step fall behind and are lost.
\item \textbf{Plain steps} (one step size per parameter group) lose most states even with exact statistics,
  because the gradient of a state carries its occupation.
\item \textbf{Gain per state} (one adaptive step size per hidden state) keeps them with exact statistics but not
  with measured ones.
\item \textbf{Visit-count rule} (\texttt{--hidden-norm occupancy}): each state keeps a leaky average of its gradient
  and a leaky count of its visits; its step is the ratio, under one occupation-weighted step scale per parameter
  group, clipped so that no parameter moves more than one learning rate per step. Schematically, per state $m$,
  \begin{equation}
    \bar g_m\leftarrow\beta\bar g_m+(1-\beta)g_m,\qquad \bar n_m\leftarrow\beta\bar n_m+(1-\beta)\hat p_m,\qquad
    \Delta\theta_m\propto-\eta\,\mathrm{clip}\Big(\frac{\bar g_m}{\bar n_m\,s},\pm1\Big),
  \end{equation}
  with $\beta=0.999$ and $s$ the group scale. Dividing by the visit count removes the occupation factor that makes
  rare states slow. It needs one counter per state and one leaky average per coupling. The visit counts can be
  taken from the same $S$ draws as the statistics (\texttt{--occ-same-draws}), so it reads no more than per-element
  Adam.
\end{itemize}

%======================================================================================================================
\section{Generation: the reverse process}
In reversed time $\tau=5-t$ the sampler integrates
\begin{equation}
  d\vx=\big(\vx-2\grad F_{\theta(\tau)}(\vx)\big)\,d\tau+\sqrt{2T}\,d\bm W_\tau ,
\end{equation}
i.e.\ the pixels move in the effective potential $2F-\tfrac12\|\vx\|^2$ at temperature $T$. Without hidden units
this is the visible model with parameters $2\theta-\theta_{\rm lin}$. With hidden units, doubling $\bm W$ and $b$ is
\emph{not} enough ($-T\log2\cosh(2a/T)\ne-2T\log2\cosh(a/T)$). An exact construction at one temperature is to
duplicate every hidden unit (two replicas at $T$ with the trained couplings), which gives $-2T\sum_j\log Z_j$
exactly; alternatives are the doubled energy at hidden temperature $2T$, or an amplified back-action force. In
simulation the drift $\vx-2\grad F$ is used directly.

%======================================================================================================================
\section{When the idealisation holds}\label{sec:physical}
The hidden units must follow their input field. Near $t=0$ the pixel relaxation time in generation is about $0.03$
(stiffness $2k-1\approx36$), and pixel noise moves $a_j$ by about $T$ in a time $T/(2\|\bm W_j\|^2)$; the hidden
correlation time must be short against both. At equal temperatures (clamped training, releases, replica
generation) detailed balance keeps the pixel marginal exactly $e^{-F/T}$ for any hidden correlation time, which
then only adds friction; with two temperatures it also heats the pixels. Bistability and speed pull against each
other: a deep barrier gives Ising-like units, a shallow one fast switching (Kramers rate $\propto e^{-\Delta U/T}$);
the way out is a much higher attempt frequency than the pixels, with barriers of a few $T$. A Potts unit built from
$M$ bistable units with finite inhibition will sometimes have zero or two units on; whether the model tolerates
that is untested.

%======================================================================================================================
\section{Evaluation}
Each generated sample is scored by three flags, calibrated on real images (training images scored
leave-one-out): \emph{diverged} (an amplitude above 15 normalised units), \emph{blend} (a 0 and a 1 overlaid: one
template of each class fits clearly better than any single template; 5\% of real images are flagged) and
\emph{far} (its nearest training image is farther than for 90\% of real images at $14\times14$; a fixed distance
at $28\times28$). A sample is \emph{clean} if it is none of these. The \emph{clean yield (50/50)} averages the clean
share over the two classes among non-diverged samples; real images score 85.7\% ($14\times14$) and 83\%
($28\times28$). Variety is measured separately by $k$-nearest-neighbour precision and recall ($k=3$) per class.
All models are sampled from the same 500 initial states and noise paths, so comparisons are paired (McNemar tests;
paired bootstrap, 95\% intervals).

%======================================================================================================================
\section{Results so far}\label{sec:results}

\begin{table}[h]\centering\small
\begin{tabular}{llr}
\toprule
problem & model & clean yield \\
\midrule
$28\times28$, 0/1 & paper's model (no hidden units) & 31.9\% \\
                  & one Ising unit with spring & 43.5\% \\
                  & Potts 128 + both, exact training & 86.3\% \\
                  & Potts 128 + both, local rule, $S=64$, jackknife & 87.6\% \\
                  & Potts 128 + both, visit-count rule, $S=16$ & 87.1\% \\
                  & Potts 512 + both, power 4, exact training & 92.3\% \\
\midrule
$14\times14$, 0/1 & Potts 128 + both, exact training & 84.3\% \\
(8 reads)         & local rule, per-element Adam, $S=8$ & 68.4\% \\
                  & local rule, visit-count rule, $S=8$, same-draw counts & 81.7\% \\
\midrule
$14\times14$, 0/1, power 4 & radius 7 / 3 / 2 & 84.7 / 83.9 / 83.8\% \\
Potts 128 + both & 8 nearest neighbours, springs only & 81.4\% \\
exact training   & 8 nearest neighbours, + optomechanical term & 81.1\% \\
                 & 4 nearest neighbours, springs only & 74.2\% \\
                 & no pair couplings & 60.0\% \\
                 & 4 nearest neighbours, springs only, no hidden units & 0\% \\
\midrule
$14\times14$, 0/1/2/3 & Potts 512 + both, exact / local ($S=64$, visit-count) & 85.0 / 74.7\% \\
\bottomrule
\end{tabular}
\caption{Main results (one training run per arm; seed-to-seed variation about one point where measured). Real
images score 83\% ($28\times28$) and 85.7\% ($14\times14$).}
\end{table}

\begin{itemize}
\item Hidden units fix the paper model's main failure (0/1 superpositions) and, with enough Potts states and the
  spring coupling, bring the samples to the level of real images; power 4 matches power 6 from about 128 states.
\item The score-matching gradient is a local three-factor rule. With measured hidden statistics it needs the
  jackknife and an optimiser that keeps rare states alive: the visit-count rule beats per-element Adam by 13
  points at equal reads (8 per configuration), replicated at a second training seed, and with 16 samples matches
  exact training at $28\times28$.
\item A much simpler visible model works: nearest-neighbour linear springs (8 neighbours) cost about 3 points; the
  optomechanical term is not needed. Without hidden units such a model produces no clean digits.
\item Open: local training on four digits ($-10$ points); quality as a function of the number of Potts states and
  of Ising units with the simplest visible model (running); several Ising units instead of one Potts unit; finite
  hidden speed and finite winner-take-all inhibition.
\end{itemize}

\end{document}
